\documentclass[11pt]{article}
\usepackage[T1]{fontenc}
\usepackage[a4paper,margin=1in]{geometry}
\usepackage{amsmath,amssymb,amsthm,mathtools}
\usepackage{booktabs,array,longtable}
\usepackage{enumitem}
\usepackage{xcolor}
\usepackage{microtype}
\usepackage[hidelinks]{hyperref}

\setlength{\parskip}{0.45em}
\setlength{\parindent}{1.4em}
\setlist[itemize]{topsep=0.25em,itemsep=0.15em,leftmargin=1.8em}
\setlist[enumerate]{topsep=0.25em,itemsep=0.15em,leftmargin=2em}

\newtheorem{theorem}{Theorem}[section]
\newtheorem{proposition}[theorem]{Proposition}
\newtheorem{lemma}[theorem]{Lemma}
\newtheorem{corollary}[theorem]{Corollary}
\theoremstyle{definition}
\newtheorem{definition}[theorem]{Definition}
\newtheorem{example}[theorem]{Example}
\newtheorem{problem}[theorem]{Problem}
\newtheorem{conjecture}[theorem]{Conjecture}
\theoremstyle{remark}
\newtheorem{remark}[theorem]{Remark}

\newcommand{\PGL}{\operatorname{PGL}}
\newcommand{\PSL}{\operatorname{PSL}}
\newcommand{\GL}{\operatorname{GL}}
\newcommand{\Aff}{\operatorname{Aff}}
\newcommand{\Biv}{\operatorname{Biv}}
\newcommand{\Hist}{\operatorname{Hist}}
\newcommand{\Stab}{\operatorname{Stab}}
\newcommand{\im}{\operatorname{im}}
\newcommand{\Id}{\operatorname{Id}}
\newcommand{\Ev}{\operatorname{Ev}}
\newcommand{\Hol}{\operatorname{Hol}}
\newcommand{\PP}{\mathbb P}
\newcommand{\F}{\mathbb F}
\newcommand{\cH}{\mathcal H}
\newcommand{\cG}{\mathcal G}
\newcommand{\cB}{\mathcal B}
\newcommand{\cC}{\mathcal C}
\newcommand{\sslash}{\mathbin{/\mkern-6mu/}}

\title{\textbf{Bilateral Arithmetic Histories and Projective Condensation}\\
\large A Bridge from Arithmetic Expression Geometry to Projective Concept--Relation Geometry}
\author{Mingli Yuan}
\date{Working note, July 2026}

\begin{document}
\maketitle

\begin{abstract}
Arithmetic Expression Geometry (AEG) starts from the principle that an
expression is not exhausted by its final value: its ordered evaluation history
is part of the object.  Existing AEG develops one constant-chirality,
affine sector of a larger sequential syntax.  This note completes that syntax
by allowing the active subexpression to occupy either input slot of each binary
operation.  The resulting marked spinal histories have a non-degenerate
projective evaluation whose image is exactly \(\PGL_2(K)\) over a field \(K\);
the familiar affine theory is the stabilizer of one point at infinity.

On the projective side, an ordered pair of distinct points of
\(\PP^1(K)\) is equivalent to a rank-one idempotent projector on \(K^2\).
Consequently, with \(G=\PGL_2(K)\) and \(H\) the stabilizer of an ordered
pair, the bivaluation space is \(G/H\).  A third projective point supplies a
full frame, hence a point of the \(H\)-torsor above the condensed result; only
after a reference lift is chosen does this acquire an \(H\)-valued coordinate.
This produces a precise process--result hierarchy:
\[
\text{marked history}
\longrightarrow
G
\longrightarrow
G/H
\longrightarrow
G/B_\pm,
\]
where each arrow forgets a different layer of process information.  The note
also distinguishes syntactic chirality from point--predicate duality, derives
the affine/Riccati double flow, gives a complete finite-field counting model,
and identifies a telescoping obstruction in the current contextual-transport
ansatz.  The final sections formulate the minimal connection and holonomy data
needed to prevent process geometry from collapsing back to endpoint geometry.
\end{abstract}

\section{Thesis, scope, and proof status}

The central thesis is that the projective layer need not be appended to AEG as
an external completion.  It can be generated internally by completing the
threadlike syntax that AEG already uses.

The note separates three levels of assertion.

\begin{enumerate}
  \item \textbf{Proved algebraic statements.}
  These include the classification of sequential expression trees, the
  generation of \(\PGL_2(K)\) by bilateral arithmetic sections, the
  identification of regular bivaluations with rank-one idempotents, the
  homogeneous-space descriptions \(G/H\) and \(G/B_\pm\), the finite-field
  counts, and the telescoping identity for an exact endpoint-defined
  transport.

  \item \textbf{Structural proposals.}
  These include the word-decorated history groupoid, group-valued comparison
  defects, and the use of the principal \(H\)-bundle \(G\to G/H\) as the
  natural home of process residue after projective condensation.

  \item \textbf{Open geometric bridges.}
  These include a canonical comparison with the contact form used in AEG,
  a projective replacement for the two-dimensional ACS area formula, and the
  semantic interpretation of rank-one condensation as a concept or predicate.
\end{enumerate}

This ordering is deliberate.  The projective algebra is independent of any
learning, probability, or epistemic interpretation.  Such interpretations
should be added only after the underlying process geometry has been made
functorial.

\section{Marked spinal expressions}

\subsection{Two one-hole sections of a binary operation}

Let \(K\) be a field and let
\[
\Omega=\{+,-,\times,/\}
\]
be the arithmetic binary signature, with the usual partial-domain convention
for division.  Given \(\omega\in\Omega\) and a constant \(c\in K\), define two
one-hole contexts
\begin{equation}
 C^{(1)}_{\omega,c}[z]:=\omega(z,c),
 \qquad
 C^{(2)}_{\omega,c}[z]:=\omega(c,z).
 \label{eq:two-sections}
\end{equation}
The superscript records the operand slot occupied by the active value.  This
notation avoids the ambiguous phrases ``left-expanded'' and
``right-expanded'', whose meanings reverse depending on whether one looks at
the recursive subtree, the newly attached leaf, or the direction of written
growth.

\begin{definition}[Marked spinal history]
A marked spinal history over \(K\) is data
\[
\gamma=
\bigl(x_0;(\omega_i,c_i,\varepsilon_i)_{i=1}^{n}\bigr),
\qquad
\varepsilon_i\in\{1,2\},
\]
with recursive states
\[
x_i=C^{(\varepsilon_i)}_{\omega_i,c_i}[x_{i-1}].
\]
The word
\[
\varepsilon(\gamma)=\varepsilon_1\varepsilon_2\cdots\varepsilon_n
\]
is its \emph{chirality word}.  The same data without the initial state
\(x_0\) is a free history.
\end{definition}

Pure slot-\(1\) and pure slot-\(2\) histories are the two combs.  Allowing an
arbitrary chirality word gives all mixed combs without introducing a branching
dependency.

\subsection{Intrinsic tree characterization}

For a binary expression tree \(T\), let \(I(T)\) be the set of its internal
vertices, ordered by dependency: \(u<v\) when the output of \(u\) is required
before \(v\) can be evaluated.

\begin{theorem}[Sequential-tree classification]
\label{thm:sequential-tree}
For a finite binary expression tree \(T\), the following are equivalent.
\begin{enumerate}
  \item The dependency poset \(I(T)\) has a unique linear extension.
  \item The dependency poset \(I(T)\) is a chain.
  \item Every internal vertex has at most one internal child.
  \item Either \(T\) is a single leaf, marked \(x\), or its internal vertices
  form a single spine.  In the latter case, after choosing either leaf of the
  innermost binary vertex as the marked accumulator \(x\), \(T\) can be
  generated by the grammar
  \[
  T::=x\mid\omega(T,c)\mid\omega(c,T).
  \]
  Here each \(c\) is one of the remaining atomic leaves.
\end{enumerate}
\end{theorem}

\begin{proof}
A finite poset has a unique linear extension only if every two elements are
comparable.  Indeed, for incomparable \(u,v\), adjoining \(u<v\) creates no
cycle and hence extends to a linear order; adjoining \(v<u\) does likewise.
The two resulting linear extensions are distinct.  Hence (1) implies (2),
and the converse is immediate.

For a binary tree, two internal children of the same vertex are incomparable
in the dependency order.  Thus (2) implies (3).  If there are no internal
vertices, mark the unique leaf.  Otherwise, if every internal vertex has at
most one internal child, following that child from the root visits all
internal vertices, so they form a spine; this gives (4).  Finally the grammar
in (4) visibly gives one internal child at each stage and hence a chain of
dependencies.
\end{proof}

\begin{remark}[Why the seed must be marked]
An unmarked comb determines a unique order of internal evaluations, but the
innermost binary node still has two leaves.  Marking the initial state specifies
which leaf is the evolving accumulator and converts the tree unambiguously
into a word of one-hole contexts.
\end{remark}

\subsection{Mirror is not reversal}

\begin{definition}[Mirror and temporal reversal]
The \emph{mirror} \(m\gamma\) flips every slot bit,
\[
1\longleftrightarrow 2,
\]
and mirrors the corresponding planar tree without changing the bottom-up time
index \(i\).  The \emph{temporal reversal} \(r\gamma\) reverses the context word.
When every step is invertible, the \emph{path inverse} additionally replaces
each context by its inverse.
\end{definition}

Mirror and reversal are independent commuting involutions at the marked-word
level.  For example,
\[
\gamma=((x+1)\times 2),\qquad
m\gamma=2\times(1+x),\qquad
r\gamma=(x\times2)+1.
\]
The first two induce \(2x+2\), whereas the reversed history induces \(2x+1\).
Thus the current affine torsion is defined by comparison with reversal/order;
it is not, without additional structure, a detector of tree mirror.

\begin{remark}[Associativity as a 2-cell]
If process is primary, associativity and commutativity should not erase marked
histories at the moment of definition.  They are better represented by
rewrite 2-cells between histories.  For one fixed infix order of atomic
leaves, the ordinary left and right combs are the extremal vertices of the
Tamari order.  Relating those extrema to the two pure \emph{marked} histories
requires a corresponding relabeling of the seed and constants, since planar
mirror alone reverses the leaf order.  Closing compatible spinal histories
under rotations leads to the full associahedron and, in general, to branching
evaluation schedules.
\end{remark}

\subsection{Bilateral resolvent recursion}

The same distinction can appear in recursive closure, but only after enough
noncommutative data have been retained.  A useful warning comes first.  For a
single fixed element \(p\) and a common central initial value
\(S_0=\widetilde S_0\) (for example \(0\) or \(1\)), the two recursions
\[
S_{n+1}=1+pS_n,
\qquad
\widetilde S_{n+1}=1+\widetilde S_n p
\]
produce the same polynomial in \(p\), even when \(p\) is a matrix or an
operator.  If \(1-p\) is invertible, their fixed-point equations
\[
S=1+pS,
\qquad
\widetilde S=1+\widetilde S p
\]
also have the same solution \((1-p)^{-1}\).  Noncommutativity alone is not
enough.

A genuine bilateral recursion needs genuinely noncommuting retained data---for
example, a seed that does not commute with \(p\), or step-dependent transports
and offsets that do not all lie in one commuting subalgebra.  The cleanest
placement test uses the same data on both sides:
\[
S_{n+1}=a_n+p_nS_n,
\qquad
\widetilde S_{n+1}=a_n+\widetilde S_np_n.
\]
Now the ordered products place new factors at opposite temporal ends, so the
two histories need not agree.  This is the appropriate operator-level test
case for bilateral process semantics before primitive division is introduced.

\section{Projective evaluation of bilateral arithmetic}

\subsection{Ordinary and projective semantics}

Let
\[
\PP^1(K)=K\cup\{\infty\}.
\]
Every non-degenerate one-hole arithmetic section in
\eqref{eq:two-sections} extends to a fractional linear transformation of
\(\PP^1(K)\).  Using the convention
\[
\begin{pmatrix}\alpha&\beta\\ \gamma&\delta\end{pmatrix}\cdot z
=\frac{\alpha z+\beta}{\gamma z+\delta},
\]
the elementary sections are represented as follows:
\[
\begin{array}{c@{\qquad}c@{\qquad}c}
\toprule
\text{section} & \text{map} & \text{projective matrix}\\
\midrule
C^{(1)}_{+,c}=C^{(2)}_{+,c}
& z\mapsto z+c
& \begin{pmatrix}1&c\\0&1\end{pmatrix}\\[3mm]
C^{(1)}_{-,c}
& z\mapsto z-c
& \begin{pmatrix}1&-c\\0&1\end{pmatrix}\\[3mm]
C^{(2)}_{-,c}
& z\mapsto c-z
& \begin{pmatrix}-1&c\\0&1\end{pmatrix}\\[3mm]
C^{(1)}_{\times,c}=C^{(2)}_{\times,c}
& z\mapsto cz
& \begin{pmatrix}c&0\\0&1\end{pmatrix}\\[3mm]
C^{(1)}_{/,c}
& z\mapsto z/c
& \begin{pmatrix}1&0\\0&c\end{pmatrix}\\[3mm]
C^{(2)}_{/,c}
& z\mapsto c/z
& \begin{pmatrix}0&c\\1&0\end{pmatrix}\\
\bottomrule
\end{array}
\]
Here \(c\neq0\) whenever invertibility requires it.  If multiplication by zero
or constant maps are admitted, the natural target is a projective matrix
monoid, not \(\PGL_2(K)\).

\begin{definition}[Non-degenerate bilateral history]
Let \(\Hist_K^{\pm,\times}\) be the free history language generated by the
invertible sections in the table above.  Projective evaluation is the monoid
map
\[
\rho:\Hist_K^{\pm,\times}\longrightarrow\PGL_2(K)
\]
obtained by multiplying the corresponding matrices in chronological
composition order.
\end{definition}

\begin{remark}[Projective continuation is not ordinary arithmetic]
The projective map \(z\mapsto c/z\) sends \(0\) to \(\infty\), whereas ordinary
arithmetic declares division by zero undefined.  A projective completion
therefore changes the semantic category.  A sound theory should retain a
domain label distinguishing a legitimate chart transition from an admissible
ordinary-arithmetic evaluation.
\end{remark}

\subsection{The bilateral closure theorem}

\begin{theorem}[Bilateral arithmetic generates \(\PGL_2\)]
\label{thm:bilateral-pgl}
Let \(K\) be a field.  The non-degenerate projective evaluations of bilateral
four-operation spinal histories generate all of \(\PGL_2(K)\).
\end{theorem}

\begin{proof}
Translations \(T_s(z)=z+s\), nonzero scalings \(D_q(z)=qz\), and the Weyl
inversion \(J(z)=-1/z\) are bilateral arithmetic sections.  It suffices to
show that they
generate every fractional linear map
\[
f(z)=\frac{Az+B}{Cz+D},
\qquad AD-BC\neq0.
\]
If \(C=0\), then \(f\) is affine and is a translation after a nonzero scaling.
If \(C\neq0\), the identity
\begin{equation}
f(z)
=
\frac{A}{C}
+
\frac{BC-AD}{C^2}\,
\frac{1}{z+D/C}
\label{eq:mobius-decomposition}
\end{equation}
equivalently gives the explicit composition
\[
f=T_{A/C}\circ D_{(AD-BC)/C^2}\circ J\circ T_{D/C}.
\]
Hence \(f\) lies in the evaluation image.
\end{proof}

\begin{corollary}[The algebraic slot-\(1\) sector is affine]
If the bilateral language is restricted to the first operand slot, the
non-degenerate four-operation sections are affine:
\[
z\mapsto az+b,\qquad a\in K^\times,\ b\in K.
\]
Their image is the Borel subgroup
\[
B_\infty=\Stab_{\PGL_2(K)}(\infty)\cong\Aff(1,K).
\]
The existing real AEG embeds in this sector; it fills the full Borel only after
its allowed multipliers are enlarged to all of \(K^\times\).
The second-slot division \(c/z\) is the elementary operation that leaves this
subgroup.
\end{corollary}

\begin{remark}[The positive real sector]
The current continuous AEG parameterizes multiplication as
\(z\mapsto e^\lambda z\) with real \(\lambda\), so over \(\mathbb R\) it sees
the orientation-preserving identity component \(\Aff^+(1,\mathbb R)\), not
the entire real Borel with arbitrary nonzero slope.  The full algebraic Borel
appears only after all multipliers in \(K^\times\) are admitted.
\end{remark}

\begin{corollary}[Bruhat interpretation]
Let \(B=B_\infty\) and let \(J(z)=-1/z\).  Then
\[
\PGL_2(K)=B\sqcup BJB
\]
at the level of Bruhat cells, and \(B\) together with \(J\) generates the
whole projective group.  Bilateral completion is therefore a syntactic
realization of the rank-one Bruhat completion of the affine sector.
\end{corollary}

\begin{example}[Finite continued fractions]
A pure slot-\(2\) nesting involving addition and division has the form
\[
a_0+\frac{c_1}{a_1+\dfrac{c_2}{a_2+\cdots}},
\]
so finite generalized continued fractions are a distinguished bilateral
threadlike sublanguage.  Their convergent matrices are precisely products in
\(\PGL_2(K)\).
\end{example}

\section{History before quotient}

\subsection{The evaluation functor}

Let \(\cH_K^\pm\) be the category whose arrows retain the full marked context
word and whose objects are admissible values (or projective points).  An arrow
\((\gamma,x):x\to\rho(\gamma)x\) is not identified with another word merely
because both words induce the same transformation.  Evaluation gives a
functor
\[
\Ev:\cH_K^\pm\longrightarrow \PP^1(K)\sslash\PGL_2(K).
\]

This yields the basic hierarchy
\begin{equation}
\text{marked syntax}
\longrightarrow
\text{one-hole context word}
\longrightarrow
\text{projective operator}
\longrightarrow
\text{endpoint value}.
\label{eq:forgetful-syntax-tower}
\end{equation}
Each arrow forgets data.  The principle ``process before result'' means that
the fibers and kernel pairs of these maps are themselves mathematical objects;
it does not mean that no quotient may ever be taken.

\begin{definition}[Relative projective defect]
For two invertible histories \(\gamma,\delta\), define
\[
\mathcal K(\gamma,\delta)
:=
\rho(\delta)^{-1}\rho(\gamma)
\in\PGL_2(K).
\]
This is a group-valued comparison with an explicit choice of reference
history.  It equals the identity exactly when the two histories induce the
same projective operator.
\end{definition}

\begin{proposition}[Affine scalar torsion as a special coordinate]
Suppose
\[
\rho(\gamma)(x)=sx+A_\gamma,\qquad
\rho(\delta)(x)=sx+A_\delta
\]
have the same nonzero linear part.  Then
\[
\rho(\delta)^{-1}\rho(\gamma)(x)
=x+\frac{A_\gamma-A_\delta}{s}.
\]
The ordinary endpoint defect
\[
\rho(\gamma)(x)-\rho(\delta)(x)=A_\gamma-A_\delta
\]
is therefore the translation coordinate of the relative holonomy, up to the
common scale \(s\).  The familiar quantity
\(\mu(e^\lambda-1)\) is the elementary two-step instance of this reduction.
\end{proposition}

\begin{remark}
Once inversion is admitted, a scalar subtraction of endpoints is no longer
the natural invariant: it can depend on the base point and encounter poles.
The group-valued defect \(\mathcal K\), or its conjugacy/stabilizer class under
the relevant endpoint constraint, survives this enlargement.
\end{remark}

\section{Left and right affine velocities}

Except for the discrete cocycle proposition below, this section works over
\(K=\mathbb R\).  The same formulas may be read over \(\mathbb C\), or
formally in characteristic zero after replacing analytic exponentials by an
appropriate Lie parameter.  They are not assertions about arbitrary fields
or finite-field differential calculus.

\subsection{Two Maurer--Cartan readings}

Write an affine element as
\[
g=
\begin{pmatrix}
e^\lambda&\xi\\
0&1
\end{pmatrix}.
\]

\begin{proposition}[Left and right affine differentials]
The translation components of the left and right Maurer--Cartan forms are
\[
\bigl(g^{-1}dg\bigr)_{\rm tr}=e^{-\lambda}d\xi,
\qquad
\bigl(dg\,g^{-1}\bigr)_{\rm tr}=d\xi-\xi\,d\lambda.
\]
\end{proposition}

\begin{proof}
This follows by multiplying
\[
g^{-1}=
\begin{pmatrix}
e^{-\lambda}&-e^{-\lambda}\xi\\0&1
\end{pmatrix}
\]
with \(dg\) on the appropriate side.
\end{proof}

The first form measures velocity in the body/source frame; the second measures
it in the spatial/target frame.  They become the differential signatures of
growth from opposite temporal ends when histories are represented by left or
right group multiplication.  This multiplication-end axis is not identical
to operand-slot chirality: a theorem connecting them requires an explicit
representation of context insertion.  If
\[
\widehat\xi:=e^{-\lambda}\xi,
\]
then
\[
d\xi-\xi\,d\lambda=e^\lambda d\widehat\xi.
\]
Thus the standard translation \(\xi\) and the source-normalized translation
\(\widehat\xi\) are not competing definitions; they are the two natural
cocycle coordinates of the same process.

\subsection{Future and past weighting}

\begin{proposition}[Two affine cocycle sums]
Let \(K\) be an arbitrary field, let
\[
f_i(x)=s_i x+t_i,\qquad s_i\in K^\times,
\]
and let \(f_n\circ\cdots\circ f_1(x)=\Phi_nx+\xi_n\).  Then
\[
\Phi_n=\prod_{i=1}^n s_i,
\qquad
\xi_n=\sum_{i=1}^n t_i\prod_{j=i+1}^n s_j,
\]
whereas the source-normalized translation is
\[
\widehat\xi_n:=\frac{\xi_n}{\Phi_n}
=\sum_{i=1}^n\frac{t_i}{\prod_{j=1}^i s_j}.
\]
\end{proposition}

\begin{proof}
The first two identities follow by induction under postcomposition by
\(f_{n+1}\).  Dividing the translation identity by
\(\Phi_n=\prod_js_j\) gives the last formula.
\end{proof}

The coordinate \(\xi_n\) weights an additive event by all future scalings,
whereas \(\widehat\xi_n\) discounts it by the accumulated past scale.  These
are the discrete counterparts of the \(e^{M}\) and \(e^{-M}\) path weights
that appear when the affine cocycles are written in charge coordinates.

\subsection{From affine flow to Riccati flow}

Return here to the real Lie setting (or its characteristic-zero formal
analogue).  Let
\[
E=\partial_z,\qquad H=z\partial_z,\qquad F=z^2\partial_z.
\]
Up to sign conventions, these close under the \(\mathfrak{sl}_2\) brackets.
The affine theory uses only \(E\) and \(H\).  Inversion supplies the missing
direction because
\[
J\,T_s\,J^{-1}(z)=\frac{z}{1-sz}
=z+s z^2+O(s^2),
\qquad J(z)=-\frac1z.
\]
Consequently the general infinitesimal projective flow is Riccati:
\begin{equation}
\dot z=\beta+\alpha z+\kappa z^2.
\label{eq:riccati-general}
\end{equation}
The current AEG flow is the Borel slice \(\kappa=0\).

\section{Bivaluation as coupled point--predicate transport}

\subsection{Projective points, copoints, and incidence}

Let \(V=K^2\) and \(V^\vee\) be its algebraic dual.  A projective point and a
projective copoint are denoted
\[
[v]\in\PP(V),\qquad [\varphi]\in\PP(V^\vee).
\]
They are incident when \(\varphi(v)=0\).  On the affine chart used in the
bivaluation notes, use the sign-reversed coordinate
\[
\chi([x:y])=-x/y.
\]
Thus \(\chi=-z\) relative to the standard coordinate of the projective
evaluation section.  Choose representatives
\[
v_a=\begin{pmatrix}-a\\1\end{pmatrix},
\qquad
\varphi_b=\begin{pmatrix}-b&1\end{pmatrix}.
\]
Their normalized local pairing is
\[
I(a,b)=\varphi_b(v_a)=1+ab.
\]
The kernel line of \(\varphi_b\) has bivaluation point coordinate
\[
p=-\frac1b.
\]
Its standard \(z\)-coordinate is \(1/b=-p\).  Thus the coordinate switch is
explicit rather than an identification of \(a\) with the earlier \(z\).
Thus \(a=p\), equivalently \(1+ab=0\), is the incidence divisor.

\begin{definition}[Regular bivaluation]
A regular bivaluation over \(K\) is an ordered pair
\[
(q^+,q^-)\in\PP^1(K)\times\PP^1(K),
\qquad q^+\neq q^-.
\]
The first line is the image/prototype direction and the second is the
kernel/predicate direction.  Write
\[
\Biv_K^{\rm reg}
:=
\PP^1(K)\times\PP^1(K)\setminus\Delta.
\]
\end{definition}

\subsection{Rank-one projector theorem}

\begin{theorem}[Bivaluations are rank-one idempotents]
\label{thm:biv-idempotent}
There is a natural bijection
\[
\Biv_K^{\rm reg}
\longleftrightarrow
\{\Pi\in\operatorname{End}(K^2):
\Pi^2=\Pi,\ \operatorname{rank}\Pi=1\}.
\]
If \(q^+=[v]\), \(q^-=\PP(\ker\varphi)\), and
\(\varphi(v)\neq0\), the corresponding projector is
\begin{equation}
\Pi_{v,\varphi}
=\frac{v\otimes\varphi}{\varphi(v)}.
\label{eq:rank-one-projector}
\end{equation}
In the local \((a,b)\)-chart this is
\[
\Pi_{a,b}
=
\frac{1}{1+ab}
\begin{pmatrix}
ab&-a\\
-b&1
\end{pmatrix}.
\]
\end{theorem}

\begin{proof}
Given distinct one-dimensional subspaces \(L^+,L^-\subset K^2\), one has
\(K^2=L^+\oplus L^-\).  There is a unique projection onto \(L^+\) along
\(L^-\), and it is a rank-one idempotent.  Choosing
\(L^+=Kv\) and a nonzero covector \(\varphi\) with
\(\ker\varphi=L^-\) gives \eqref{eq:rank-one-projector}.

Conversely, if \(\Pi^2=\Pi\), then
\[
K^2=\im\Pi\oplus\ker\Pi.
\]
Rank one makes both summands lines, and they are distinct.  The two
constructions are inverse.
\end{proof}

\subsection{Covariant and contravariant histories}

For \(G\in\PGL(V)\), point and predicate representatives transform by
\[
v\longmapsto Gv,
\qquad
\varphi\longmapsto\varphi G^{-1}.
\]
Hence the unnormalized pairing is invariant:
\[
(\varphi G^{-1})(Gv)=\varphi(v),
\]
and the projector transforms by conjugation,
\[
\Pi_{v,\varphi}\longmapsto G\Pi_{v,\varphi}G^{-1}.
\]

\begin{proposition}[Local affine bivaluation law]
Suppose the point coordinate transforms by
\[
a\longmapsto Aa+B,\qquad A\in K^\times.
\]
In the standard \(z=-a\) coordinate this is represented by
\[
G_{A,B}=\begin{pmatrix}A&-B\\0&1\end{pmatrix},
\qquad z\longmapsto Az-B.
\]
Then the predicate zero and normalized covector coordinate transform as
\begin{equation}
p\longmapsto Ap+B,
\qquad
b\longmapsto\frac{b}{A-Bb},
\qquad p=-\frac1b.
\label{eq:bivaluation-affine-law}
\end{equation}
For infinitesimal addition and scaling over \(\mathbb R\) (or in a
characteristic-zero formal differential setting),
\[
s=\mu\,du,\qquad k=e^{\lambda\,dv},
\]
the three local flows are
\begin{equation}
\begin{aligned}
da&=\mu\,du+\lambda a\,dv,\\
db&=\mu b^2\,du-\lambda b\,dv,\\
dp&=\mu\,du+\lambda p\,dv.
\end{aligned}
\label{eq:bivaluation-double-flow}
\end{equation}
\end{proposition}

\begin{proof}
The point \(p\) is a projective point and therefore obeys the same affine law
as \(a\).  Substituting \(p=-1/b\) gives
\[
b'=-\frac1{Ap+B}=\frac{b}{A-Bb}.
\]
Expanding the addition update
\[
b\longmapsto\frac{b}{1-\mu b\,du}
\]
and the scaling update \(b\mapsto e^{-\lambda\,dv}b\) to first order gives
\eqref{eq:bivaluation-double-flow}.
\end{proof}

\begin{remark}[Two distinct binary axes]
Operand-slot chirality and point--predicate duality must not be identified.
The former belongs to syntax:
\[
\omega(z,c)\quad\hbox{versus}\quad\omega(c,z).
\]
The latter belongs to a representation and its contragredient:
\[
v\mapsto Gv\quad\hbox{versus}\quad\varphi\mapsto\varphi G^{-1}.
\]
They meet naturally inside \(\PGL_2\), but a theorem relating mirror to
transpose, inverse, or dual transport requires an explicit representation and
an anti-involution.  It is not automatic.
\end{remark}

\section{Projective frames and the process--result bundle}

\subsection{Homogeneous-space descriptions}

Fix a reference regular pair \((q_0^+,q_0^-)\) and set
\[
G=\PGL_2(K),
\qquad
H=\Stab_G(q_0^+,q_0^-).
\]
Also write \(B_+=\Stab_G(q_0^+)\) and
\(B_-=\Stab_G(q_0^-)\).  Then \(H=B_+\cap B_-\), and each \(B_\pm\) is a
conjugate affine Borel subgroup.

\begin{theorem}[Projective quotient tower]
\label{thm:quotient-tower}
There are natural identifications
\[
G/H\cong\Biv_K^{\rm reg},
\qquad
G/B_+\cong\PP^1(K),
\qquad
G/B_-\cong\PP^1(K).
\]
The maps
\[
G/H\longrightarrow G/B_+,
\qquad
G/H\longrightarrow G/B_-
\]
forget respectively the kernel or image member of the ordered pair.
\end{theorem}

\begin{proof}
The action of \(\PGL_2(K)\) on ordered pairs of distinct projective points is
transitive.  The stabilizer of the reference pair is exactly \(H\), so the
orbit--stabilizer identification gives \(G/H\).  The same argument for a
single point gives the two Borel quotients.  The forgetful maps are induced by
the inclusions \(H\subset B_\pm\).
\end{proof}

\subsection{A third point restores a full frame}

\begin{definition}[Framed bivaluation]
A framed regular bivaluation is a pairwise distinct ordered triple
\[
(q^+,q^-,q^\infty)\in\bigl(\PP^1(K)\bigr)^3.
\]
The first two entries determine the projector result; the third specifies a
projective arithmetic frame or chart.
\end{definition}

\begin{theorem}[Frame torsor]
\label{thm:frame-torsor}
After choosing one reference triple, the set of framed regular bivaluations is
a \(G\)-torsor.  The projection
\[
(q^+,q^-,q^\infty)\longmapsto(q^+,q^-)
\]
has fibers on which \(H\) acts freely and transitively.  Equivalently,
\[
G\longrightarrow G/H
\]
is the process-frame bundle over the bivaluation result space.
\end{theorem}

\begin{proof}
A projective transformation is uniquely determined by the images of three
distinct points, so the \(G\)-action on ordered projective frames is simply
transitive.  For a fixed ordered pair, its stabilizer is \(H\).  In coordinates
where the pair is \((0,\infty)\), \(H\) consists of
\[
z\longmapsto\lambda z,\qquad\lambda\in K^\times,
\]
which acts freely and transitively on
\(\PP^1(K)\setminus\{0,\infty\}=K^\times\).
\end{proof}

Over an arbitrary field, ``bundle'' here means the algebraic or set-theoretic
\(H\)-torsor supplied by the quotient.  An Ehresmann connection on it belongs
to the real/complex Lie-group realization and is additional structure.

Combining syntax, operator, and condensation gives the precise forgetful
tower
\begin{equation}
\boxed{
\Hist_K^{\pm,\times}
\xrightarrow{\ \rho\ }
G
\xrightarrow{\ \pi_H\ }
G/H\cong\Biv_K^{\rm reg}
\xrightarrow{\ \pi_+\ }
G/B_+\cong\PP^1(K).
}
\label{eq:main-forgetful-tower}
\end{equation}
The algebraic slot-\(1\) completion is not the final quotient in this display;
it is the Borel subgroup \(B_\infty=\Stab(q^\infty)\subset G\) obtained by
holding one arithmetic frame point fixed.  The current positive-real AEG is
the identity-component sector inside that Borel.

\begin{proposition}[Same result, stabilizer-valued residue]
Let
\[
\pi_H(g)=g\Pi_0g^{-1}
\]
for a reference rank-one projector \(\Pi_0\).  Then
\[
\pi_H(g_1)=\pi_H(g_2)
\quad\Longleftrightarrow\quad
g_2^{-1}g_1\in H.
\]
Thus two operator-level processes with the same condensed projector differ by
an \(H\)-valued frame residue.
\end{proposition}

\begin{proof}
Equality of the conjugated projectors is equivalent to
\[
(g_2^{-1}g_1)\Pi_0(g_2^{-1}g_1)^{-1}=\Pi_0,
\]
which says precisely that \(g_2^{-1}g_1\) stabilizes the ordered image and
kernel lines of \(\Pi_0\).
\end{proof}

\begin{remark}[Where extension data can live]
A static idempotent already splits the vector space:
\[
V=\im\Pi\oplus\ker\Pi.
\]
It therefore carries no nontrivial linear extension class by itself.  Process
information discarded by condensation must instead live in such objects as
\[
\begin{gathered}
\ker\rho,\qquad \text{the \(H\)-lift in }G\to G/H,\\
\text{a connection and its holonomy},\qquad
\text{subdominant spectral data}.
\end{gathered}
\]
This is a sharper location for ``implicit'' or ``unsaid'' structure than the
projector alone.
\end{remark}

\section{The finite-field calibration}

\subsection{Exact counts}

Let \(K=\F_q\).  Then
\[
|\PP^1(\F_q)|=q+1,
\qquad
|\Biv_{\F_q}^{\rm reg}|=(q+1)q.
\]
The projective group has order
\[
|\PGL_2(\F_q)|=q(q^2-1)=(q+1)q(q-1).
\]
For a regular ordered pair,
\[
H\cong\F_q^\times,
\qquad |H|=q-1.
\]
These identities are not merely compatible counts: they are exactly the
fiber cardinalities of
\[
\PGL_2(\F_q)\longrightarrow
\Biv_{\F_q}^{\rm reg}.
\]

\begin{example}[\(\F_7\)]
For \(q=7\),
\[
|\PGL_2(\F_7)|=336,
\qquad
|\Biv_{\F_7}^{\rm reg}|=56,
\qquad
|H|=6.
\]
Equivalently, there are \(8\cdot7\cdot6=336\) ordered projective frames,
\(8\cdot7=56\) ordered image--kernel pairs, and six possible frame points
above each projector result.  This is a fully enumerable model of process
frame versus condensed result.
\end{example}

\subsection{Three failures that must not be conflated}

\begin{enumerate}
  \item \textbf{Incidence failure.}
  If \(q^+=q^-\), or locally \(1+ab=0\), the projector
  \eqref{eq:rank-one-projector} is undefined.  This is an intrinsic
  degeneration of the transverse image--kernel splitting.  The unnormalized
  rank-one operator \(v\otimes\varphi\) still exists there, but it satisfies
  \((v\otimes\varphi)^2=0\) rather than defining a nonzero idempotent.

  \item \textbf{Chart escape.}
  A denominator such as \(A-Bb\) may vanish while the global ordered pair and
  projector remain regular.  The correct response is a chart change, not a
  declaration that the projective object has ceased to exist.

  \item \textbf{Failure of attracting dynamics.}
  A rank-one projector may exist algebraically even when no sequence of
  projective transformations converges to it in the intended topology.
\end{enumerate}

\begin{proposition}[No automatic condensation over a finite field]
The existence of a rank-one idempotent over \(\F_q\) does not imply dynamical
condensation.  Powers of an element of \(\PGL_2(\F_q)\) are periodic, and a
sequence in the discrete finite group converges only if it is eventually
constant.
\end{proposition}

Thus algebraic condensation
\[
\Pi^2=\Pi,\qquad\operatorname{rank}\Pi=1
\]
must be distinguished from an analytic or probabilistic limit
\[
\frac{\widetilde G_n\cdots \widetilde G_1}{\rho_n}
\longrightarrow\Pi
\quad\text{in }\operatorname{End}(K^2),
\]
where compatible lifts \(\widetilde G_i\in\GL_2(K)\) have first been chosen.
Such a singular limit does not occur inside \(\PGL_2(K)\) itself, and requires
additional norm, topology, valuation, positivity, or stochastic structure.

\subsection{Split, parabolic, and nonsplit process types}

The conjugacy type of a projective operator gives a first finite failure and
dynamics atlas:
\begin{itemize}
  \item the identity fixes all \(q+1\) projective points;
  \item nonidentity split semisimple elements have two fixed points over
  \(\F_q\);
  \item parabolic/unipotent elements have one repeated fixed point;
  \item nonsplit semisimple elements have no fixed point over \(\F_q\), but
  acquire two over a quadratic extension;
  \item all orbits remain finite and periodic.
\end{itemize}
The nonsplit case is especially instructive: absence of a pole in the base
field is not absence of projective pole data.  It is an extension-field
visibility phenomenon.

\section{An exact-transport obstruction in the current relation ansatz}

\subsection{The telescoping identity}

Work here in the real positive-affine group, where the logarithmic scale
\(\eta\in\mathbb R\) is single-valued.  Consider lifted vertices
\((a_t,\eta_t)\) and define an affine frame
\[
h_t(z)=e^{\eta_t}z+a_t.
\]
We write \((b,\lambda)_{\rm aff}\) for the map
\(z\mapsto e^\lambda z+b\).
The endpoint-anchored relation used in the current relation-theory draft is
\begin{equation}
r_t(z)
=e^{\eta_t-\eta_{t-1}}z
+a_t-e^{\eta_t-\eta_{t-1}}a_{t-1}.
\label{eq:exact-edge}
\end{equation}

\begin{proposition}[Exact edge labels telescope]
\label{prop:exact-telescope}
The edge in \eqref{eq:exact-edge} is
\[
r_t=h_t h_{t-1}^{-1}.
\]
Consequently,
\begin{equation}
r_t r_{t-1}\cdots r_{s+1}
=h_t h_s^{-1}
=
\left(
a_t-e^{\eta_t-\eta_s}a_s,\,
\eta_t-\eta_s
\right)_{\rm aff},
\label{eq:telescope}
\end{equation}
which is independent of all intermediate lifted vertices.
\end{proposition}

\begin{proof}
The first identity follows by direct substitution.  The product then cancels
pairwise:
\[
(h_t h_{t-1}^{-1})(h_{t-1}h_{t-2}^{-1})
\cdots(h_{s+1}h_s^{-1})=h_t h_s^{-1}.
\]
\end{proof}

\begin{corollary}
For this exact ansatz, a closed lifted path has trivial holonomy.  A path whose
base value closes but whose lift does not close records only the endpoint
frame difference \(\eta_t-\eta_s\), not the intermediate path.
\end{corollary}

This does not invalidate the use of endpoint relations.  It shows that this
\emph{globally exact}, single-frame ansatz defines a pure-gauge transport and
therefore cannot by itself realize the desired principle that context residue
depends on intermediate history.  It is not a claim that every flat local
connection has trivial global monodromy: transition functions, multivalued
frames, or a nontrivial bundle can retain it.

\subsection{Minimal repairs}

There are four compatible ways to retain process.

\begin{enumerate}
  \item \textbf{Word-decorated path groupoid.}
  Keep the free marked history category and the evaluation functor
  \[
  \cH^\pm\xrightarrow{\Ev}\PP^1\sslash\PGL_2.
  \]
  Distinct histories remain distinct even when their evaluated morphisms
  coincide.

  \item \textbf{Independent edge connection.}
  Do not force every edge label to be the exact coboundary
  \(h_t h_{t-1}^{-1}\).  For example, write
  \[
  r_t=h_t k_t h_{t-1}^{-1}
  \]
  with an independently specified or learned transport \(k_t\).  If
  \(a_t=h_tq_0\) and the prescribed point endpoints must be preserved, require
  \(k_t\in\Stab_G(q_0)\); for a prescribed bivaluation path, require
  \(k_t\in H\).  Without that stabilizer condition the construction changes
  the base path as well as its transport.  Closed products can then retain the
  conjugacy class of the \(k_t\)-product.

  \item \textbf{Connection on \(G\to G/H\).}
  Regard a bivaluation path as a path in \(G/H\), choose horizontal lifts in
  \(G\), and measure the \(H\)-valued mismatch of a closed lift.  This is the
  natural differential-geometric realization of the stabilizer residue for
  real or complex Lie groups.  Over a general field, \(G\to G/H\) should first
  be treated as an algebraic \(H\)-torsor; a connection requires additional
  geometric structure.

  \item \textbf{Rewrite 2-cells.}
  Treat associativity, commutativity, distributivity, and admissible program
  rewrites as 2-cells.  Their costs and coherence loops live above operator
  equality rather than being erased by it.
\end{enumerate}

\begin{problem}[Projective torsion]
Construct a projective comparison invariant that:
\begin{enumerate}
  \item reduces to the affine translation defect in the Borel sector;
  \item is equivariant under change of projective frame;
  \item records an \(H\)-valued holonomy when two histories have the same
  bivaluation result; and
  \item admits a surface or cocycle formula analogous to the ACS weighted-area
  identity.
\end{enumerate}
\end{problem}

\section{Hyperbolic and contact geometry: a precise test}

Over \(K=\mathbb R\), the orientation-preserving projective group satisfies
\[
\PSL_2(\mathbb R)/\mathrm{SO}(2)\cong\mathbb H^2.
\]
Moreover, after choosing a reference oriented unit tangent vector,
\[
\PSL_2(\mathbb R)\cong T^1\mathbb H^2,
\]
and the unit tangent bundle carries its canonical contact form.  This gives a
structural reason that a two-dimensional hyperbolic model and a
three-dimensional contact model can arise from the same projective completion.

\begin{conjecture}[Contact comparison]
There is a Gauss or Iwasawa coordinate chart on an open cell of
\(\PSL_2(\mathbb R)\) in which the canonical contact distribution on
\(T^1\mathbb H^2\) is related, by an explicit bundle map and conformal factor,
to the AEG horizontal distribution
\[
\ker\bigl(da-\mu\,du-\lambda a\,dv\bigr).
\]
\end{conjecture}

Local contactomorphism alone would not establish this claim, because Darboux's
theorem makes all three-dimensional contact structures locally equivalent.
The required comparison must also intertwine:
\begin{itemize}
  \item the arithmetic generators with projective flows;
  \item the affine bracket with the appropriate Maurer--Cartan equation;
  \item the projection to the hyperbolic base;
  \item and the group-valued holonomy of finite histories.
\end{itemize}

\section{A minimal theorem and experiment program}

\subsection{The first rigorous package}

The following package can be completed without invoking learning semantics.

\begin{enumerate}
  \item Formalize marked spinal expressions and prove the sequential-tree
  classification.
  \item Define ordinary partial evaluation and projective evaluation as
  distinct functors.
  \item Prove the bilateral closure theorem and its Bruhat refinement over an
  arbitrary field.
  \item Prove
  \[
  \Biv_K^{\rm reg}\cong G/H
  \]
  and the frame-torsor theorem.
  \item Define relative group-valued defects at the history, operator,
  point-result, and projector-result levels.
  \item Prove the finite-field counts and generate complete tables for small
  \(q\).
  \item Replace the exact contextual transport by either a word-decorated
  groupoid or a non-flat connection and exhibit the first nontrivial loop.
\end{enumerate}

\subsection{Finite falsification tests}

Small finite fields provide exact unit tests:
\begin{itemize}
  \item enumerate \(\Hist^\pm\) up to a word-length cutoff and compare its
  operator fibers under \(\rho\);
  \item verify the surjection onto \(\PGL_2(\F_q)\);
  \item enumerate all rank-one idempotents and their \(H\)-fibers;
  \item classify history images as split, parabolic, or nonsplit;
  \item lift nonsplit fixed points to \(\F_{q^2}\);
  \item distinguish incidence failure from chart escape in every local chart.
\end{itemize}
These tests are finite, reproducible, and independent of analytic convergence.

\subsection{Only then add semantic condensation}

The interpretation
\[
\text{projector}=\text{concept}=\text{prototype}\otimes\text{predicate}
\]
is a substantive hypothesis, not a consequence of rank-one linear algebra.
It becomes informative only if point--predicate models outperform appropriate
point-only controls or explain new transport/holonomy observables.  The
projective mathematics should therefore remain usable even if a particular
learning hypothesis is inconclusive or refuted.

\section{Relation to the two working projects}

This note was motivated by two existing bodies of work.

\begin{itemize}
  \item In the AEG repository, the early syntax-tree material already displays
  both pure combs and mixed combs, while the later affine theory develops only
  one constant-chirality sector.  The Möbius grid \(z\mapsto-1/z\), the ACS
  torsion formula, contact calculus, and loop/equality program supply the
  local ingredients used here.

  \item In the PCRG paper repository, the APG draft explicitly asks for a
  projective completion of history-sensitive arithmetic; the bivaluation notes
  identify
  \[
  \Biv_K^{\rm reg}
  =\PP^1_K\times\PP^1_K\setminus\Delta
  \]
  with rank-one projectors; and the finite-field notes distinguish incidence,
  chart, extension-field, and dynamical residues.  The present note supplies
  a candidate arithmetic syntax for that projective completion and sharpens
  the location of process residue.
\end{itemize}

The proposed division of labor is:
\[
\boxed{
\begin{array}{c}
\text{bilateral AEG explains where projective transport comes from},\\[1mm]
\text{APG/PCRG explains what projective histories can condense into},\\[1mm]
\text{the history and frame fibers explain what condensation forgets}.
\end{array}}
\]

\section{Conclusion}

The missing opposite-slot thread is not merely another family of examples.
It changes the natural evaluation group from the affine stabilizer of one
projective point to the full rank-one projective group.  Once this completion
is made, the point--predicate bivaluation space appears as a homogeneous
quotient, and a full projective frame becomes a lift of a condensed
image--kernel result.

The resulting hierarchy
\[
\Hist_K^{\pm,\times}
\longrightarrow
\PGL_2(K)
\longrightarrow
\Biv_K^{\rm reg}
\longrightarrow
\PP^1(K)
\]
is a concrete mathematical realization of the statement that process carries
more information than result.  Its fibers are not philosophical decoration:
they are word relations, stabilizer directions, chart choices, and potential
holonomies.  A viable bridge between AEG and PCRG should be built on these
fibers before it is built on semantic analogy.

\end{document}
